\PassOptionsToPackage{dvipsnames,svgnames,table}{xcolor}
\documentclass[11pt]{article}
\usepackage[a4paper,margin=22mm]{geometry}
\usepackage[no-math]{fontspec}
\setmainfont{DejaVu Serif}
\usepackage{amsmath,amssymb,amsthm,mathtools,graphicx,xcolor,hyperref,xurl,xspace,ifthen,enumerate}
\usepackage[shortlabels]{enumitem}
\setlength{\emergencystretch}{4em}
\AtBeginDocument{\special{pdf:mapline pzdr Dingbats <uzdr.pfb}}
\SetMathAlphabet{\mathbf}{normal}{OT1}{cmr}{bx}{n}
\SetMathAlphabet{\mathbf}{bold}{OT1}{cmr}{bx}{n}
\newtheorem{theo}{Theorem}[section]
\newtheorem{lemm}[theo]{Lemma}
\newtheorem{thm}[theo]{Theorem}
\newtheorem{proposition}[theo]{Proposition}
\newtheorem{corollary}[theo]{Corollary}
\newtheorem{definition}[theo]{Definition}
\newtheorem{theorem}[theo]{Theorem}
\newtheorem*{remas*}{Remarks}
\newtheorem{remas}[theo]{Remarks}
\newtheorem*{exam*}{Example}
\newtheorem{ques}[theo]{Question}
\providecommand{\og}{«\nobreakspace}
\providecommand{\fg}{\nobreakspace»}
\newtheorem{lem}[theo]{Lemma}
\newtheorem{lemma}[theo]{Lemma}
\newtheorem{prop}[theo]{Proposition}
\newtheorem{coro}[theo]{Corollary}
\newtheorem{corol}[theo]{Corollary}
\newtheorem*{conj*}{Conjecture}
\newtheorem{conj}[theo]{Conjecture}
\newtheorem{conjecture}[theo]{Conjecture}
\newtheorem{Proposition}[theo]{Proposition}
\newtheorem{theoreme}[theo]{Theorem}
\theoremstyle{definition}
\newtheorem{defi}[theo]{Definition}
\newtheorem{exam}[theo]{Example}
\newtheorem{exem}[theo]{Example}
\newtheorem{remark}[theo]{Remark}
\newtheorem{example}[theo]{Example}
\newtheorem{exams}[theo]{Examples}
\newtheorem{rema}[theo]{Remark}
\newtheorem{nota}[theo]{Notation}
\newtheorem{notation}[theo]{Notation}
\newtheorem*{notas}{Notations}
\newtheorem{result}[theo]{Result}
\newtheorem{question}[theo]{Question}
\newtheorem{prob}[theo]{Problem}
\newtheorem*{rema*}{Remark}
\newtheorem*{defi*}{Definition}
\newtheorem*{theo*}{Theorem}
\newtheorem*{lemm*}{Lemma}
\newtheorem*{prop*}{Proposition}
\newcommand{\enoncename}{}
\newtheorem{corpusenonce}[theo]{\enoncename}
\newenvironment{enonce}[2][]{\renewcommand{\enoncename}{#2}\begin{corpusenonce}}{\end{corpusenonce}}
\newenvironment{enonce*}[2][]{\par\medskip\noindent\textbf{#2.}\itshape}{\par\medskip}
\newenvironment{proofc}[1][\proofname]{\begin{proof}[#1]}{\end{proof}}
\newcommand{\Mc}{\mathcal{M}}
\newcommand{\pointir}{.\ ---\ }
\newcommand{\equalenv}[2]{\expandafter\let\csname#1\expandafter\endcsname\csname#2\endcsname\expandafter\let\csname end#1\expandafter\endcsname\csname end#2\endcsname}
\providecommand{\diagbox}[3][]{\begin{tikzpicture}[x=1em,y=1em,baseline=(current bounding box.center)]\useasboundingbox(0,0) rectangle(3,1.4);\draw(0,1.4)--(3,0);\node at(.5,.3){#2};\node at(2.5,1.1){#3};\end{tikzpicture}}
\providecommand{\eg}{e.g.\ }
\providecommand{\ie}{i.e.\ }
\providecommand{\cf}{cf.\ }
\providecommand{\longthanks}[1]{\par\smallskip\noindent #1\par\smallskip}
\providecommand{\qbinom}[2]{\genfrac{[}{]}{0pt}{}{#1}{#2}}

\providecommand{\mathof}[1]{\mathrm{#1}}

\numberwithin{equation}{section}


\def\endappendix{}


\usepackage{latexsym,amsmath,amssymb,verbatim,tikz,tikz-cd}
\usepackage{mathtools}
\usepackage{graphicx,xcolor}

\usepackage[english]{babel}
\usepackage{amssymb}
\usepackage{amsmath}
\usepackage{amsthm}
\usepackage{graphicx}
\usepackage{csquotes}
\usepackage{bm}
\usepackage{enumitem}
\usepackage{subcaption}
\usepackage{caption}
\renewcommand{\thesubfigure}{\Alph{subfigure}}
\usepackage{dirtytalk}
\providecommand{\nopunct}{\spacefactor1007\relax}
\newtheorem{problem}[theo]{Problem}

\newcommand{\sym}{\mathrm{Sym}}
\newcommand{\qsym}{\mathrm{QSym}}
\newcommand{\Des}{\operatorname{Des}}
\newcommand{\Asc}{\operatorname{Asc}}
\newcommand{\Syt}{\operatorname{SYT}}
\newcommand{\Short}{\operatorname{Short}}
\newcommand{\Loop}{\operatorname{Loop}}
\newcommand{\Stable}{\operatorname{Stable}}
\newcommand{\len}{\operatorname{len}}
\newcommand{\height}{\operatorname{height}}
\newcommand{\core}{\operatorname{core}}
\newcommand{\res}{\operatorname{res}}
\newcommand{\row}{\operatorname{row}}
\newcommand{\insrt}{\operatorname{insert}}
\newcommand{\UP}{\operatorname{UP}}
\newcommand{\DOWN}{\operatorname{DOWN}}


\newcommand{\ZZ}{\mathbb{Z}}
\newcommand{\NN}{\mathbb{N}}
\newcommand{\RR}{\mathbb{R}}
\newcommand{\QQ}{\mathbb{Q}}
\newcommand{\PP}{\mathbb{P}}

\newcommand{\F}{\mathcal{F}}
\newcommand{\Q}{\mathcal{Q}}
\newcommand{\A}{\mathcal{A}}
\newcommand{\M}{\mathcal{M}}
\newcommand{\I}{\mathcal{I}}
\newcommand{\PM}{\mathcal{PM}}
\newcommand{\Pcal}{\mathcal{P}}


\newcommand{\odds}{\mathrm{odds}}

\makeatletter
\newcommand{\proofstep}[1]{
  \par
  \addvspace{\medskipamount}
  \textbf{#1\@addpunct{.}}\enspace\ignorespaces
}

\newcommand\restr[2]{{
  \left.\kern-\nulldelimiterspace 
  #1 
  \vphantom{\big|} 
  \right|_{#2} 
  }}
\makeatother












\title{The map sending a matching to its core and its standard Young tableau is well defined}
\author{Avichai Marmor}
\date{doi:10.5802/alco.351; CC BY 4.0}
\begin{document}
\maketitle
\noindent\textit{Editorial note.} The counted statement is the author's unlabelled corollary environment at publisher-source.tex lines 854-860; its proof environment opens at line 861, immediately after \textbackslash{}end\{corollary\}, so the pairing is by adjacency and the retained proof range 861-869 is the whole environment.\par\medskip
\noindent\textit{Editorial note.} The proposition prop:core well defined, statement and proof (source lines 846-853), is retained as context because the counted proof cites it at source line 862; it is the only outside label of the retained proof.\par\medskip
\noindent\textit{Editorial note.} Because the corollary carries no author label, its identity is recorded by the pinned source position of the two environments (statement lines 854-860, proof lines 861-869) and by the claim key, which encodes that position.\par\medskip
\noindent\textit{Editorial note.} This article carries no label-number dump and the retained counter values are left empty for the whole source, so the printed numbering of the card is the wrapper's default; the original printed number of this unlabelled corollary is not recoverable from the pinned material and is therefore not asserted.\par\medskip
\noindent\textit{Editorial scope.} One original result of the pinned publisher article, counted once: well-definedness of the map from a matching to its core and its two-row tableau (unlabelled corollary at publisher-source.tex lines 854-860).\par\medskip
\noindent\textit{Allowed standard prerequisites (editorial).}\begin{itemize}
\item The union of cores and two-row tableaux. The target is the disjoint union over k of the matchings with short-chord-free core and the standard Young tableaux of shape (N-k,k).. It is the codomain in which the map F is asserted to be well defined.
\item Core and two-row tableau of a matching. The core is the matching on the stable vertices and T(m) is the two-row tableau whose second row records the unstable chords.. They are the two components of the map.
\item Standard Young tableau. A two-row tableau of shape (N-k,k) is standard when both rows and both columns are increasing and each letter occurs once.. It is the property the proof verifies for T(m).
\end{itemize}\medskip
\begin{proposition}
\label{prop:core well defined}
    Given a matching, its stable vertices and its core are well-defined.
\end{proposition}
\begin{proof}
    Let $m \in \M_N$ be a matching. Denote the chords that are removed during a given reduction process of $m$ by $e_1, \dots, e_k$, and denote $e_\ell = (i_\ell,j_\ell)$. Moreover, denote the chords that are removed during another reduction process by $e_1', \dots, e_{k'}'$, and denote $e_\ell' = (i_\ell',j_\ell')$. It suffices to prove that $e_\ell \in \{e_1', \dots, e_{k'}'\}$ for all $1 \le \ell \le k$ (\ie every chord that is removed during the first reduction process is removed during the second process as well). Assume by contradiction that $e_\ell \notin \{e_1', \dots, e_{k'}'\}$ for some $\ell$, and denote by $\ell_0$ the minimal such $\ell$.
    That is, $e_1,\dots,e_{\ell_0-1} \in \{e_1', \dots, e_{k'}'\}$. Removing the chords $e_1, \dots, e_k$ from $m$ constitutes a valid reduction process, so the chord $e_{\ell_0}$ becomes short before it is removed. That is, $i \in \{i_1,j_1,\dots,i_{\ell_0-1},j_{\ell_0-1}\}$ for all $i_{\ell_0} < i < j_{\ell_0}$. Therefore, the chord $e_{\ell_0}$ is a short chord of the core obtained by removing $e_1', \dots, e_{k'}'$ of $m$, contradicting the requirement that the reduction process continues until there are no short chords remaining. Therefore, we may conclude that if a chord is removed during a reduction process then it is removed during any reduction process, and $\core(m)$ and $\Stable(m)$ are well-defined.
\end{proof}

\begin{corollary}
    Let $N$ and $f$ be nonnegative integers. Then the function
    \[
        F:\M_{N,f} \to \bigcup_{k=0}^n \M_{N-2k,f}(\Short = \emptyset) \times \Syt(N-k,k)
    \]
    defined by $F(m) = (\core(m),T(m))$ is well-defined.
\end{corollary}

\begin{proof}
    Let $m\in \M_{N,f}$ be a matching. By Proposition~\ref{prop:core well defined}, we obtain that $\core(m)$ and $\Stable(m)$ are well-defined. Therefore, the two-row tableau $T:=T(m)$ consisting of $N$ cells that is defined by
    \[
        \row_2(T) = \{j \mid \text{the chord } (i,j) \in m \text{ is unstable}\}
    \]
    is also well-defined. In order to prove that $F$ is well-defined, it remains to show that $T(m) \in \Syt(N-k,k)$, where $k$ is the number of unstable chords of $m$.

    Obviously, every letter $i \in [N]$ appears exactly once in $T(m)$, and the rows are increasing. Notice that for every entry $j \in \row_2(T(m))$ there exists $i<j$ such that $(i,j) \in m$ is an unstable chord. Therefore, every entry $j \in \row_2(T(m))$ is associated to an entry $i \in \row_1(T(m))$ such that $i<j$. Moreover, if $j\ne j' \in \row_2(T(m))$ and $(i,j),\; (i',j') \in m$ then $i \ne i'$, so the columns of $T(m)$ are increasing and $T(m)$ is a standard Young tableau. Finally, $|\row_2(T(m))| = k$, implying that $T(m) \in \Syt(N-k,k)$.
\end{proof}

\end{document}
